\chapter{Конструкторская часть}
%<пишем про алгоритмы, архитектуру ПО и тп>

\section{Описание алгоритмов}

Схемы последовательного и распараллеленного алгоритмов представлены на рис. 2.1 и 2.2 соответственно.

\begin{figure}[!htbp]
	\centering
	\includesvg[width=0.2\textwidth]{lab4_single}
	\caption{Схема последовательного алгоритма}
\end{figure}
\FloatBarrier

\begin{figure}[!htbp]
	\centering
	\includesvg[width=0.3\textwidth]{lab4_par}
	\caption{Схема параллельного алгоритма}
\end{figure}
\FloatBarrier

\section*{Вывод}
%<что сделали в конструкторке и получили в результате, кратко>

В данном разделе мной были построены схемы последовательного и параллельного алгоритмов. % и вычислена их трудоемкость с помощью введенной модели вычислений.

\clearpage
